\section{Wprowadzenie}
Celem projektu jest detekcja choroby Chagasa z EKG pacjentów, zgodnie z George B. Moody
Physionet Challenge 2025: \url{https://moody-challenge.physionet.org/2025/}

Choroba Chagasa jest tropikalną chorobą pasożytniczą występującą w Ameryce Łacińskiej.
Jej objawy możemy podzielić na dwie fazy:
\begin{itemize}
    \item fazę ostrą - podobną do grypy; w niektórych przypadkach dodatkowo występuje guzek
    o nazwie chagoma, powiększenie węzłów chłonnych lub obrzęk w okolicy oczodołu,
    \item fazę przewlekłą - przeważnie bezobjawową; w ok. 20-30\% przypadków objawia się
    m.in. chorobą mięśnia sercowego.
\end{itemize}

Zgodnie z opisem autorów challengu, symptomy choroby Chagasa można rozpoznać w EKG.